\subsection{3. Die sechste Dezimalstelle ist richtig - aber nur für die
eine zusätzliche
Richtung}\label{die-sechste-dezimalstelle-ist-richtig---aber-nur-fuxfcr-die-eine-zusuxe4tzliche-richtung}

Die neue numerische Diagonalisierung ergibt:

{\def\LTcaptype{none} % do not increment counter
\begin{longtable}[]{@{}
  >{\raggedright\arraybackslash}p{(\linewidth - 6\tabcolsep) * \real{0.2000}}
  >{\raggedleft\arraybackslash}p{(\linewidth - 6\tabcolsep) * \real{0.2667}}
  >{\raggedleft\arraybackslash}p{(\linewidth - 6\tabcolsep) * \real{0.2667}}
  >{\raggedleft\arraybackslash}p{(\linewidth - 6\tabcolsep) * \real{0.2667}}@{}}
\toprule\noalign{}
\begin{minipage}[b]{\linewidth}\raggedright
g/Δ
\end{minipage} & \begin{minipage}[b]{\linewidth}\raggedleft
vierdimensionale Ritz-Matrix
\end{minipage} & \begin{minipage}[b]{\linewidth}\raggedleft
plus w₂
\end{minipage} & \begin{minipage}[b]{\linewidth}\raggedleft
Änderung
\end{minipage} \\
\midrule\noalign{}
\endhead
\bottomrule\noalign{}
\endlastfoot
1/20 & −1,0942308394409612 & −1,0942318710142458 & −1,0315733·10⁻⁶ \\
1/40 & −0,29535718178032033 & −0,29535720496683315 & −2,3186513·10⁻⁸ \\
1/100 & −0,04789424798601009 & −0,04789424801311775 &
−2,7107663·10⁻¹¹ \\
\end{longtable}
}

Das sind numerische Eigenwerte exakter kleiner Ritz-Matrizen, keine nach
außen gerundeten vollständigen Spektraleinschließungen. Durch
Eliminieren der einzigen neuen w₂-Koordinate bekommt die v₃-Diagonale
die exakt bestimmte energieabhängige Korrektur

\[\adjustbox{max width=\linewidth}{$\displaystyle 
-\frac{g^2(1809/47632)}{2\Delta-E}.
$}\]

Das erklärt die kleine Änderung durch diesen einen Zweig. Es sagt nichts
über die ausgelassenen höheren Bosonlagen.

\subsubsection{Eigene einfache Folgeherleitung: Das vollständige
Ritz-Residuum}\label{eigene-einfache-folgeherleitung-das-vollstuxe4ndige-ritz-residuum}

Sei ψ₃ der normierte Ritz-Eigenvektor in \codeword{span(v₀,…,v₃)}, und
c₃ seine Komponente entlang \codeword{v₃/√ν₃}. Außerhalb dieses Raums
führt Hψ₃ genau in zwei orthogonale Richtungen: v₄ und w₂. Deshalb gilt
exakt

\[\adjustbox{max width=\linewidth}{$\displaystyle 
\|(H-E_{\rm Ritz})\psi_3\|^2
=g^2|c_3|^2\frac{\nu_4+\|w_2\|^2}{\nu_3}.
$}\]

Hier wird \codeword{ν₄=952296652800} aus dem früher vollständig
geprüften Momentbeleg übernommen; die vierte Ordnung wurde in dieser
Runde nicht erneut enumeriert.

Bei g/Δ=1/20 folgt numerisch ein Residuum von rund \textbf{0,373063 Δ}.
Der w₂-Zweig trägt weniger als \textbf{23 Millionstel} des quadrierten
Residuums. Die große ausgelassene Richtung ist v₄, nicht w₂. Aus dem
winzigen Nichtschlussanteil darf daher keine praktisch vollständige
Konvergenz abgeleitet werden.

Außerdem liegt die neue Fünf-Zustands-Ritzenergie mindestens
\textbf{0,0354 Δ über} der bereits bewiesenen Obergrenze des wahren
Grundzustands. Das ist eine nachweislich noch relevante Lücke,
unabhängig von einer Konvergenzschätzung.

\subsection{4. Der bereits bewiesene native Grundsatz bleibt
gültig}\label{der-bereits-bewiesene-native-grundsatz-bleibt-guxfcltig}

Der ältere, gepinnte Beweis lautet: Für Δ\textgreater0 und
0\textless\textbar g\textbar/Δ≤1/20 besitzt der unveränderte
μ=0-Hamiltonoperator einen eindeutigen globalen Grundzustand in N=64; er
ist ein Spin(10)×SU(4)-Singulett. Bei 1/20 bestehen außerdem positive
Lücke und die bekannten Energie-, Besetzungs- und Polschranken.

Die neue Ausführung benutzt schwächere Sektor-Untergrenzen und weniger
Ritz-Vektoren. Dass diese schwächere Methode bei 1/20 nicht genügt, ist
weder eine Widerlegung noch ein Wiederöffnen des früheren Satzes.
Insbesondere sind die vorgeschlagenen Ziele \codeword{E<−1,2Δ} bzw.
\codeword{E<−1,1625Δ} unmöglich, weil bereits \codeword{E₀>−1,158089Δ}
bewiesen ist. Bessere Schranken für die Konkurrenzsektoren sind
erforderlich, nicht eine unrealistisch tiefere Energie.

Die neuen Zahlen \codeword{Nb≈0,84211266}, \codeword{Zh≈0,97368398} und
\codeword{εh≈0,03107309Δ} gehören zum nichtstationären K₃-Ritz-Zustand.
Nb liegt sogar knapp unter der früher bewiesenen Schranke
\codeword{Nb(Ω)>0,842846}. Zh bezeichnet die gesamte Entnahmenorm dieser
Näherung, nicht das Gewicht der isolierten niedrigen Linie. Ein erstes
Moment auf einem Ritz-Zustand ist keine neu bestimmte Polenergie auf Ω.

Offen bleiben der vollständige Ω-Vektor, die physische Auswahl des
μ=0-Vertrags und die Herkunft ausführbarer räumlicher Operationen.

\subsection{5. Feldwörterbuch: bestätigte Algebra und notwendige
Einschränkungen}\label{feldwuxf6rterbuch-bestuxe4tigte-algebra-und-notwendige-einschruxe4nkungen}

Frisch exakt geprüft wurden:

\begin{itemize}
\tightlist
\item
  Der Trägergraph ist zusammenhängend, 15-regulär und dreiecksfrei.
\item
  Der Fünf-Zyklus \codeword{16→45→0→30→37→16} liegt wirklich im
  Trägergraphen.
\item
  Für jeden der 60 nativen Kanäle verschwindet der gleichhändige
  Ein-Kopie-Weyl-Skalaransatz \codeword{M\_A⊗ε} als Grassmann-Bilinear.
\item
  Bei Dirac-Kernen sind C, Cγ⁵ und alle vier Cγ\^{}μγ⁵ antisymmetrisch;
  alle vier Cγ\^{}μ und alle sechs Cσ\^{}\{μν\} sind symmetrisch.
\item
  Die 43 getesteten nichttrivialen Gradierungen ergeben genau die
  angegebenen Kanalzahlen 24/36, 40/20 bzw. 32/28. Die
  Stabilisatordimensionen 36, 52 und 28 wurden diesmal durch exakt
  diagonale reelle Gram-Matrizen der Kommutatorbedingungen bestätigt,
  nicht durch einen SVD-Schwellwert.
\end{itemize}

Der Fünf-Zyklus allein schließt nur eine diagonale
±-Händigkeitszuweisung aus, die an jedem Paar entgegengesetzte
Händigkeiten verlangt. Er allein ist kein basisunabhängiger Ausschluss
und kein Ausschluss zusätzlicher Feldkomponenten.

\subsubsection{Eigener stärkerer Satz: auch ein Basiswechsel repariert
den Ein-Kopie-Vektoransatz
nicht}\label{eigener-stuxe4rkerer-satz-auch-ein-basiswechsel-repariert-den-ein-kopie-vektoransatz-nicht}

Sei Γ eine hermitesche Händigkeit auf dem ursprünglichen
64-dimensionalen internen Raum. Für einen ausschließlich
entgegengesetzt-händigen, unveränderten Paartensor müsste

\[\adjustbox{max width=\linewidth}{$\displaystyle 
\Gamma^T M_A+M_A\Gamma=0\quad\text{für alle }A
$}\]

gelten. Adjungieren und Einsetzen zeigt

\[\adjustbox{max width=\linewidth}{$\displaystyle 
[\Gamma,M_A^\dagger M_B]=0\quad\text{für alle }A,B.
$}\]

Die native Faktorisierung wurde exakt rekonstruiert:

\[\adjustbox{max width=\linewidth}{$\displaystyle 
M_{k,a}=S_k\otimes C_a,\quad
\sum_k S_k^\dagger S_k=5I_{16},\quad
\sum_a C_a^\dagger C_a=3I_4.
$}\]

Daher muss Γ sowohl mit allen \codeword{S\_k†S\_l⊗I₄} als auch mit allen
\codeword{I₁₆⊗C\_a†C\_b} kommutieren. Die von den ersten Produktmatrizen
erzeugte Algebra ist die volle M₁₆; die zweite ist M₄. Hierfür enthält
die Prüfung vollständige Rangzertifikate von ganzzahligen Matrixwörtern
modulo 101: \textbf{256 unabhängige Spinorwörter bis Länge zwei und 16
unabhängige Farbprodukte bis Länge eins}. Voller Rang modulo einer
Primzahl beweist vollen Rang über den komplexen Zahlen; es wird
ausdrücklich keine Rangdefizienz von einem endlichen Körper übertragen.

Folglich ist Γ skalar. Die ursprüngliche Gleichung erzwingt dann Γ=0.
Insbesondere gibt es \textbf{keine hermitesche Involution Γ²=I} mit der
verlangten Eigenschaft. Diese Schlusskette schließt den
Ein-Weyl-pro-Label-Vektoradapter jetzt unabhängig von einer Wahl der
internen Basis aus.

Der Satz betrifft unveränderte W-Kanäle und unveränderten
64-dimensionalen Träger. Er schließt \textbf{nicht} Diracfelder pro
Label, zusätzliche Kopien, Ableitungsadapter oder andere ausdrücklich
erweiterte Modelle aus.

\subsection{6. Die minimale Erweiterung darf nicht aus Versehen mit
ausgeschlossen
werden}\label{die-minimale-erweiterung-darf-nicht-aus-versehen-mit-ausgeschlossen-werden}

Zwei explizite Gegenkontrollen wurden auf allen 60 Kanälen neu
gerechnet:

\begin{enumerate}
\def\labelenumi{\arabic{enumi}.}
\tightlist
\item
  \textbf{Vierkomponenten-Diracfelder pro Label:} \codeword{M\_A⊗Cγ⁰}
  ist nichtverschwindend und antisymmetrisch im gesamten
  Grassmann-Index. Es enthält zusätzliche linke und rechte
  Freiheitsgrade pro ursprünglichem Label. Deshalb widerspricht es dem
  Ein-Kopie-Satz nicht. Eine chirale Standardmodell-Entstehung folgt
  daraus nicht.
\item
  \textbf{Unabhängige Zweier-Erweiterung:}
  \codeword{M\_A⊗ε\_aux⊗ε\_Lorentz} ist nichtverschwindend und
  antisymmetrisch. Der skalare Kanal wird möglich, weil die zusätzlich
  antisymmetrische interne Zweierform die Symmetrie des internen
  Paarfaktors umkehrt. Das ist die bereits bekannte minimale Reparatur
  im deklarierten Tensorprodukt-Ansatz, keine aus der Quelle
  hergeleitete Verdopplung.
\end{enumerate}

Eine solche Zweier-Erweiterung macht auch den Trägergraphen durch
\codeword{M\_A⊗σ\_x} bipartit: \codeword{Γ=I₆₄⊗diag(1,−1)}
antikommutiert exakt mit allen so erweiterten Matrizen. Damit ist
sichtbar, welche zusätzliche Ressource den ursprünglichen Ausschluss
überwindet.

``Der Vermittler kann nie Skalar sein'' muss folglich auf den
unveränderten Ein-Kopie-Ansatz begrenzt werden. Ebenso fixiert Schur nur
die interne Matrixstruktur innerhalb eines irreduziblen Blocks; es
beweist nicht die Einzigartigkeit aller möglichen Lorentz- und
Ableitungsterme. Eine gesunde Kinetik, Nebenbedingungen, Eichherkunft
und ein dynamischer Spin-2-Sektor sind hier nicht konstruiert.

\subsection{Reproduktion und Status}\label{reproduktion-und-status}

\codeword{contracted\_krylov.cpp} ist der unabhängige ganzzahlige
Kontraktionskern. \codeword{verify.py} prüft diesen Kern, die
korrigierte Pythagoras-Rechnung, kleine Ritz-Matrizen, sämtliche
angeführten Graph-, Kern- und Gradierungsidentitäten sowie die
Matrixalgebra-Zertifikate. Die Ergebnisse stehen in
\codeword{audit.json} und \codeword{audit\_optimized.json}.

Es bestehen \textbf{376 Prüfbedingungen: 371 exakte und fünf
numerische}. Für eine eigenständige Reproduktion zuerst den kleinen
C++-Kern mit C++17 und Optimierung übersetzen; anschließend
\codeword{verify.py} normal und mit \codeword{-OO} starten. Sämtliche
Eingaben liegen eingefroren unter \codeword{sources/}; die
Ergebnisdateien enthalten ihre SHA-256-Werte. Fremde Pythonprogramme
werden hierbei nicht ausgeführt. Die numerischen Ritz-Prüfungen sind als
solche markiert, alle Algebra-Zertifikate und die vollständige direkte
Kontraktion verwenden exakte Ganzzahl-/Bruchrechnung.

Beide abschließenden Ausführungen bestanden und lieferten byteidentische
JSON-Dateien mit SHA-256
\codeword{148047028236a887de1b28a0a9c12d84037a7c633892abc51dd9047f0b7a7baf}.

Keine T1-T8-Abschlussbehauptung. Keine Änderungen an fremden Quellen,
Hauptpaper, Webseite, Ledger oder Git-Historie.

\begin{center}\rule{0.5\linewidth}{0.5pt}\end{center}

\section{Anhang C: minimaler Transfer und seine
Quellen-Grenze}\label{anhang-c-minimaler-transfer-und-seine-quellen-grenze}

\section{Minimaler Ursprungsstrang: geteilte Moden und ein echter
nativer
Vierzustandskanal}\label{minimaler-ursprungsstrang-geteilte-moden-und-ein-echter-nativer-vierzustandskanal}

Stand: 15. September 2026. Unabhängiger, begrenzter Forschungsstrang.

\subsection{Ergebnis und Reichweite}\label{ergebnis-und-reichweite}

Es gibt zwei verschiedene Ergebnisse, die ausdrücklich nicht verwechselt
werden dürfen:

\begin{enumerate}
\def\labelenumi{\arabic{enumi}.}
\tightlist
\item
  \textbf{Ein exakt lösbarer Überlappungsbaustein:} Eine gemeinsame
  Fermionmode und ein gemeinsamer Bosonkanal werden in jedem festen
  positiven Ressourcensektor zu einer gewöhnlichen effektiven
  Fermionmode. Dafür ist keine direkte Endpunkt-Hopping-Wechselwirkung
  nötig. Dieser Dreierstern ist jedoch \textbf{kein Ausschnitt einer
  tatsächlichen nativen W-Zeile}.
\item
  \textbf{Ein tatsächlicher W-basierter Vierzustandskanal:} Im
  unveränderten 64-Fermion/60-Boson-Paartensor existiert ein exakt
  invariantes Vierzustandsystem, wenn ein \textbf{zusätzlicher reiner
  Bosonmischer} zwischen Kanal 0 und 1 zugelassen und ein bestimmter
  Zustand mit Gesamtladung 9 präpariert wird. Der vollständige native
  Hamiltonoperator plus dieser Mischer überträgt dort eine innere
  Fermionmarke mit einer streng abgesicherten Wahrscheinlichkeit von
  \textbf{mehr als 99,3 \%}.
\end{enumerate}

Das zweite Ergebnis benötigt keine Projektion, die die übrigen nativen
Zustände künstlich wegschneidet: Der Vierzustandsraum ist unter allen 60
ursprünglichen Paaroperatoren plus dem Mischer exakt invariant. Es ist
aber \textbf{weder eine Untersuchung der Entnahmeantwort des
N=64-Grundzustands noch ein Transport zwischen zwei unabhängigen
Banken}. Die Herkunft des Bosonmischers, die spezielle Präparation und
die räumliche Bedeutung der Marken bleiben offen.

\subsection{1. Warum eine Clockphase die bisherige Paritätsschranke
nicht
aufhebt}\label{warum-eine-clockphase-die-bisherige-parituxe4tsschranke-nicht-aufhebt}

Seien zwei operational unabhängig definierte Banken mit lokalen
Paritäten

\[\adjustbox{max width=\linewidth}{$\displaystyle 
\Pi_x=(-1)^{N_{f,x}},\qquad \Pi_y=(-1)^{N_{f,y}}
$}\]

gegeben. Kommutiert jede ausführbare lokale Operation und jeder einzelne
realisierte Messzweig mit beiden Paritäten, gilt dies auch für beliebige
Produkte, Summen, Adjungierte und starke beschränkte Grenzwerte. Eine
Clockoperation, die lediglich innerhalb jeder Bank Fermionmoden
zahlenerhaltend permutiert, bleibt in dieser Algebra.

Ein echter Einfermiontransfer erfüllt dagegen

\[\adjustbox{max width=\linewidth}{$\displaystyle 
\Pi_x f_y^\dagger f_x=-f_y^\dagger f_x\Pi_x,
\qquad
\Pi_y f_y^\dagger f_x=-f_y^\dagger f_x\Pi_y.
$}\]

Er kann daher nicht allein durch mehr Kompositionen jener Operationen
entstehen. Eine bloß paritätskovariante CP-Abbildung ist nicht die
gleiche Voraussetzung: Sie kann ungerade Krauszweige besitzen. Die
Aussage gilt nur für den ausdrücklich geraden ausführbaren
Operationssatz.

Bei geteilten Fermionmoden sind die beiden lokalen CAR-Algebren von
Anfang an nicht unabhängig. Liegt dieselbe Mode s in beiden, kann sie
nicht zugleich als zwei unabhängige antikommutierende Kopien behandelt
werden: \(\{s,s^\dagger\}=1\). Auch lokale Ladungen, die s doppelt
mitzählen, sind keine unabhängigen additiven Bankladungen. Überlappung
widerlegt somit den Paritätssatz nicht, sondern ändert seine
Voraussetzungen.

\subsection{2. Exakter CAR-Baustein aus einem geteilten Boson und
Fermion}\label{exakter-car-baustein-aus-einem-geteilten-boson-und-fermion}

Betrachtet werden drei CAR-Moden \(f_x,s,f_y\), ein CCR-Boson b und

\[\adjustbox{max width=\linewidth}{$\displaystyle 
Q_x=b^\dagger s f_x,\qquad Q_y=b^\dagger s f_y,
\qquad
H_\star=\Delta N_b+g(Q_x+Q_y+Q_x^\dagger+Q_y^\dagger).
$}\]

Der Ressourcenzähler

\[\adjustbox{max width=\linewidth}{$\displaystyle 
K=N_b+n_s
$}\]

kommutiert mit allen diesen Operationen. Auf jedem exakten K-Sektor mit
ganzzahligem \(K\ge1\) setze

\[\adjustbox{max width=\linewidth}{$\displaystyle 
c^\dagger=\frac{b^\dagger s}{\sqrt K},\qquad
c=\frac{s^\dagger b}{\sqrt K}.
$}\]

Direkt aus CCR und CAR folgen

\[\adjustbox{max width=\linewidth}{$\displaystyle 
(c^\dagger)^2=c^2=0,\qquad
\{c,c^\dagger\}=\frac{N_b+n_s}{K}=1.
$}\]

Die Mode c antikommutiert mit beiden Endpunktmoden. Ferner

\[\adjustbox{max width=\linewidth}{$\displaystyle 
N_b=K-1+n_c,
\qquad
N_f+2N_b=2K-1+n_x+n_c+n_y.
$}\]

Damit gilt exakt

\[\adjustbox{max width=\linewidth}{$\displaystyle 
\boxed{
H_\star=\Delta(K-1)+\Delta n_c
 +g\sqrt K\big(c^\dagger f_x+c^\dagger f_y+\mathrm{h.c.}\big).
}
$}\]

Die gewöhnliche Drei-Moden-Transferkette wurde hier aus einer
überlappenden Paarumwandlung erhalten. Ein nützlicher Operatorausdruck
derselben Tatsache ist

\[\adjustbox{max width=\linewidth}{$\displaystyle 
\boxed{[Q_y^\dagger,Q_x]=K f_y^\dagger f_x.}
$}\]

Das ist zunächst eine Operatoridentität. Die Verfügbarkeit getrennter
Kommutator-Kontrollsequenzen folgt daraus nicht automatisch. Schon der
konstante Summen-Hamiltonoperator zeigt jedoch Transfer. Der minimale
positive Ressourcensektor K=1 benötigt am Anfang eine besetzte
gemeinsame Fermionmode s und ein leeres Boson.

